\chapterhead{61}{ORR-22}{Solvency, Liquidity, and Other Regulation After the GFC}{4}{0}

\lohead{61}{a}{Describe and calculate the stressed VaR introduced in Basel 2.5 and
calculate the market risk capital charge.}

\orsub{Basel II and the 2007--2009 financial crisis: market risk in focus}

\begin{lettered}
\item Basel II's Advanced Internal Ratings-Based (IRB) approach allowed banks to
use their own estimates for credit risk (PD, LGD, EAD).
\item Critics argued this led to underestimation of risks, contributing to the
crisis.
\end{lettered}

\begin{orkeybox}
{\sffamily\bfseries\fontsize{8.6}{10}\selectfont\color{navy}The response: Basel
2.5 and market risk}\par
\vspace{2pt}
Introduced in 2011, Basel 2.5 aimed to improve market risk capital requirements,
and requires banks to calculate a VaR using data from a 250-day period of stressed
market conditions, known as \textbf{Stressed Value at Risk (SVaR)}.
\end{orkeybox}

\begin{itemize}
\item \textbf{Incremental Risk Charge (IRC):} a new requirement aimed at capturing
the risk from non-trading book items that could affect the trading book.
\item \textbf{Comprehensive Risk Measure (CRM):} applies to instruments sensitive
to correlations between default risks, such as securitisations and related
instruments. Ensures that banks account for complex interactions between multiple
financial instruments.
\end{itemize}

\begin{orexambox}{Example: total market risk capital charge}
Spartan State Bank has calculated a market risk VaR for the previous day equal to
\$15.6 million. The average VaR over the last 60 days is \$4.8 million. The bank
has calculated a stressed VaR for the previous day equal to \$17.7 million and an
average stressed VaR equal to \$18.4 million. Spartan State Bank has an accurate
risk measurement model and recorded only two exceptions while backtesting actual
losses against the calculated VaR. As such, the multiplicative factors, both $m_r$
and $m_s$, are set to 3. \textbf{Calculate} the total market risk capital charge.

\vspace{6pt}
\noindent{\sffamily\bfseries\fontsize{8.6}{10}\selectfont\color{forest}Solution}\par
\[
\text{total capital charge}=\$15.6\text{ million}+(\$18.4\times3)=\$70.8\text{ million}
\]
\end{orexambox}

% ---------------------------------------------------------------------
\lohead{61}{b}{Explain the process of calculating the incremental risk capital
charge for positions held in a bank's trading book.}

\orsub{Incremental Default Risk Charge (IDRC)}

\begin{lettered}
\item Before the crisis of 2007--09, trading book capital requirements were lower
than banking book requirements. This incentivised banks to move risky assets to
the trading book.
\item So the Incremental Default Risk Charge (IDRC) introduced a 99.9\% VaR,
one-year horizon for default-sensitive trading book instruments, which would make
banking book capital equal to trading book capital.
\end{lettered}

\begin{ornotebox}{Shortcomings of the IDRC and the 2007--2009 crisis}
Losses were largely due to credit downgrades, spread widening, and liquidity
problems, not just defaults. So the revised \textbf{Incremental Risk Charge (IRC)}
expanded from default risk to broader credit risks, including downgrades.
\end{ornotebox}

\orsub{Key IRC elements}

\begin{enumerate}
\item[\textcolor{rust}{\sffamily\bfseries 1.}] \textbf{Liquidity horizons.}
  \begin{checks}
  \item Banks estimate the time to sell each instrument (e.g., 6 months for a
  bond).
  \item The assumption is the bank replaces downgraded or defaulted instruments
  within this timeframe.
  \end{checks}
\item[\textcolor{rust}{\sffamily\bfseries 2.}] \textbf{Constant level of risk
assumption.}
  \begin{checks}
  \item Assumes the bank maintains a similar risk profile by replacing downgraded
  instruments.
  \end{checks}
\item[\textcolor{rust}{\sffamily\bfseries 3.}] \textbf{Specific Risk Charge (SRC).}
  \begin{checks}
  \item Considers changes in credit spreads due to downgrades.
  \end{checks}
\end{enumerate}

% ---------------------------------------------------------------------
\lohead{61}{c}{Describe the comprehensive risk (CR) capital charge for portfolios
of positions that are sensitive to correlations between default risks.}

\orsub{Comprehensive Risk Measure (CRM)}

\begin{lettered}
\item The CRM is a single capital charge replacing the incremental risk charge and
the specific risk charge for instruments dependent on credit correlation.
\item CRM is a very important measure for a portfolio of instruments such as
asset-backed securities (ABSs) and collateralised debt obligations (CDOs) that are
sensitive to the correlation between the default risks of different assets.
\item For unrated instruments or instruments rated below BB$-$, the bank must carry
dollar-for-dollar capital to back the position, equivalent to a 100\% capital
charge.
\end{lettered}

% ---------------------------------------------------------------------
\lohead{61}{d}{Define in the context of Basel III and calculate where appropriate:
Tier 1 capital and its components; Tier 2 capital and its components; required
Tier 1 equity capital, total Tier 1 capital, and total capital.}

\orsub{Basel III capital requirements}

\begin{lettered}
\item Elimination of Tier 3 capital.
\item Focus on tangible Tier 1 capital.
\item Systemic risk identification: categorisation and stricter rules for
systemically important financial institutions (SIFI).
\item Introduction of leverage ratios.
\item Liquidity requirements: new standards to ensure banks maintain adequate
liquid reserves.
\item Counterparty credit risk management.
\item Limits on exposures to single counterparties to reduce systemic risk.
\end{lettered}

\begin{center}
\begin{tikzpicture}[font=\fontsize{8}{9.8}\selectfont,
  hd/.style={rounded corners=3pt,inner sep=5pt,text=white,align=center,
             font=\sffamily\bfseries\fontsize{8.4}{10}\selectfont},
  bd/.style={text width=8.0cm,align=left,rounded corners=3pt,inner sep=7pt,
             draw=palenavy,line width=0.6pt,text=navy},
  pc/.style={font=\sffamily\bfseries\fontsize{9.6}{11.4}\selectfont}]

\node[hd,fill=navy,text width=8.0cm,anchor=north west] (h1) at (-7.2,0)
  {Common Equity Tier 1 (CET1)};
\node[bd,fill=paleblue,anchor=north west] (b1) at (-7.2,-0.56) {
  1.~Common equity $+$ retained earnings\\[2pt]
  2.~Unrealised gain/loss $+$ minority interest};
\node[pc,text=navy,anchor=west] at (1.35,-0.95) {4.5\%};

\node[hd,fill=teal,text width=8.0cm,anchor=north west] (h2) at (-7.2,-2.0)
  {Additional Tier 1 (AT1)};
\node[bd,fill=paleblue,anchor=north west] (b2) at (-7.2,-2.56) {
  1.~Non-cumulative perpetual preferred stock\\[2pt]
  2.~Convertible debt (CoCo)\\[2pt]
  3.~Other minority interest};
\node[pc,text=teal,anchor=west] at (1.35,-3.25) {1.5\%};

\draw[rust,line width=1.4pt] (2.85,-0.15) -- (3.45,-0.15) -- (3.45,-4.25) -- (2.85,-4.25);
\node[pc,text=rust,anchor=west,align=left] at (3.65,-2.2) {Total Tier 1\\ 6\% RWA};

\node[hd,fill=forest,text width=8.0cm,anchor=north west] (h3) at (-7.2,-4.95)
  {Tier 2 Capital};
\node[bd,fill=palegreen,anchor=north west] (b3) at (-7.2,-5.51) {
  1.~Subordinated debt $>5$ years\\[2pt]
  2.~Preferred stock\\[2pt]
  3.~General loan loss reserves};
\node[pc,text=forest,anchor=west] at (1.35,-6.2) {2\% RWA};
\end{tikzpicture}
\end{center}
\orfigcap{The Basel III capital stack, with the minimum ratios for each layer.}

\begin{ornotebox}{Remember}
Total capital (Tier 1 $+$ Tier 2) $=$ 8\% of RWA.\par
\vspace{5pt}
Capital is adjusted \emph{downward} for: (a) defined benefit pension plan deficits,
and (b) certain cross-holdings within a group.
\end{ornotebox}

% ---------------------------------------------------------------------
\lohead{61}{e}{Describe the motivations for and calculate the capital conservation
buffer and the countercyclical buffer, including special rules for globally
systemically important banks (G-SIBs).}

\orsub{Capital conservation buffer and countercyclical buffer}

The capital conservation buffer (CCB) is meant to protect banks in times of
financial distress.

\begin{center}
\begin{tikzpicture}[font=\fontsize{8.4}{10.4}\selectfont,
  bx/.style={text width=4.4cm,align=center,rounded corners=3pt,inner sep=8pt,
             text=navy,draw=navy,line width=0.8pt}]
\node[bx,fill=paleblue] (a) at (-5.0,0) {\textbf{7\%}\\ Core Tier 1 equity};
\node[bx,fill=paleblue] (b) at ( 0.0,0) {\textbf{8.5\%}\\ total Tier 1 capital};
\node[bx,fill=paleblue] (c) at ( 5.0,0) {\textbf{10.5\%}\\ total capital};
\node[text=rust,font=\sffamily\bfseries\fontsize{8.6}{10.4}\selectfont]
  at (0,-1.35) {in NORMAL TIMES ONLY};
\end{tikzpicture}
\end{center}
\orfigcap{Minimum ratios once the capital conservation buffer is included.}

Additionally, global systemically important banks (G-SIBs) must hold larger
buffers, ranging from 1\% to 3.5\%, to reduce their failure risk due to their
significant impact on the global financial system.

\orsubb{Countercyclical buffer (CCyB)}

To protect against cyclicality of banks. Varying from 0\% to 2.5\% of RWAs, it is
set based on the economic cycle to mitigate credit bubbles and financial
instability, with requirements differing across countries.

% ---------------------------------------------------------------------
\lohead{61}{f}{Describe and calculate ratios intended to improve the management of
liquidity risk, including the required leverage ratio, the liquidity coverage
ratio, and the net stable funding ratio.}

\orsub{Liquidity risk management}

Basel III improvements for liquidity risk management:

\begin{enumerate}
\item[\textcolor{rust}{\sffamily\bfseries 1)}] \textbf{Leverage ratio} (capital /
total exposure) of \textbf{3\%}. Total exposure is all balance sheet items plus
off-balance sheet items like loan commitments.
\item[\textcolor{rust}{\sffamily\bfseries 2)}] \textbf{Liquidity Coverage Ratio
(LCR).} Ensures banks can cover net cash outflows for 30 days during stress.
\[
\mathrm{LCR}=\frac{\text{high-quality liquid assets}}
                  {\text{net cash outflows in a 30-day period}}\ \geq 100\%
\]
\item[\textcolor{rust}{\sffamily\bfseries 3)}] \textbf{Net Stable Funding Ratio
(NSFR).} Ensures banks maintain stable funding over a one-year period. It compares
the available stable funding against the required stable funding, each adjusted by
stability factors.
\[
\mathrm{NSFR}=\frac{\text{amount of available stable funding}}
                   {\text{amount of required stable funding}}\ \geq 100\%
\]
\end{enumerate}

\begin{orexambox}{Example: LCR and NSFR}
Bank of the Bluegrass has the following balance sheet.

\vspace{5pt}
\noindent\begin{minipage}{\textwidth}
\renewcommand{\arraystretch}{1.2}
\begin{tabularx}{\textwidth}{@{}L{0.30\textwidth}L{0.10\textwidth}L{0.34\textwidth}L{0.10\textwidth}@{}}
\rowcolor{navy} \thd{Assets} & \thd{} & \thd{Liabilities and equity} & \thd{}\\
{\tabfont Cash (coins and banknotes)} & {\tabfont 10} & {\tabfont Retail deposits (less stable)} & {\tabfont 100}\\
{\tabfont Central bank reserves} & {\tabfont 10} & {\tabfont Wholesale deposits (20\% mature in one month)} & {\tabfont 75}\\
{\tabfont Treasury bonds ($>1$ yr)} & {\tabfont 10} & {\tabfont Tier 2 capital} & {\tabfont 2}\\
{\tabfont Mortgages} & {\tabfont 30} & {\tabfont Tier 1 capital} & {\tabfont 18}\\
{\tabfont Retail loans ($<1$ yr)} & {\tabfont 30} & & \\
{\tabfont Small business loans ($<1$ yr)} & {\tabfont 90} & & \\
{\tabfont Fixed assets} & {\tabfont 15} & & \\
\rowcolor{palenavy}
{\tabfont \textbf{Total assets}} & {\tabfont \textbf{195}} & {\tabfont \textbf{Total liabilities and equity}} & {\tabfont \textbf{195}}\\
\end{tabularx}
\end{minipage}

\vspace{7pt}
\begin{enumerate}
\item[\textcolor{forest}{\sffamily\bfseries 1.}] Use the balance sheet to
\textbf{calculate} the liquidity coverage ratio. Assume HQLA factors (i.e.,
haircuts) are 100\% for cash, central bank reserves and Treasury bonds (meaning no
haircut), and 0\% for all loans and fixed assets (meaning they do not count at all
in HQLA, there is a 100\% haircut and a 0\% factor). Assume also a 5\% run-off rate
for less stable retail deposits and a 100\% run-off rate of wholesale deposits that
mature in the next 30 days. Assume a 0\% run-off of Tier 2 and Tier 1 capital.
\item[\textcolor{forest}{\sffamily\bfseries 2.}] Using the ASF and RSF factors,
\textbf{calculate} the bank's net stable funding ratio.
\end{enumerate}

\vspace{7pt}
\noindent{\sffamily\bfseries\fontsize{8.6}{10}\selectfont\color{forest}Solution}\par
\vspace{3pt}
\textbf{1.} $\mathrm{HQLA}=1.0\times(30)+(0.0\times165)=30$
\[
\text{Net cash outflow in a 30-day period}=(100\times0.05)+(75\times0.2)+(20\times0.0)=20
\]
\[
\mathrm{LCR}=30/20=1.50=150\%
\]

\vspace{4pt}
\textbf{2.} $\mathrm{ASF}=(100\times0.9)+(75\times0.5)+(2\times1.0)+(18\times1.0)=\$147.50$
\[
\mathrm{RSF}=(10\times0)+(10\times0)+(10\times0.05)+(30\times0.65)+(30\times0.85)+(90\times0.85)+(15\times1.0)=\$137.00
\]
\[
\mathrm{NSFR}=147.50/137.00=1.0766=107.66\%
\]
\end{orexambox}

% ---------------------------------------------------------------------
\lohead{61}{g}{Describe the mechanics of contingent convertible bonds (CoCos) and
explain the motivations for banks to issue them.}

\orsub{Contingent convertible bonds}

\begin{ordefbox}{CoCos}
Contingent convertible bonds (CoCos) automatically convert from debt to equity
under specific financial stress conditions.
\end{ordefbox}

CoCos serve as a financial tool for banks to maintain debt in normal times,
preserving return on equity (ROE), while providing a safety net of additional
equity during financial crises.

\orsubb{Triggers for conversion include}

\begin{enumerate}
\item[\textcolor{rust}{\sffamily\bfseries 1.}] A drop in the bank's Tier 1 equity
capital ratio.
\item[\textcolor{rust}{\sffamily\bfseries 2.}] Banking supervisors determine that
the bank requires public sector support to avoid insolvency.
\item[\textcolor{rust}{\sffamily\bfseries 3.}] A specified minimum ratio of the
bank's market capitalisation to its assets is not met.
\end{enumerate}

% ---------------------------------------------------------------------
\lohead{61}{h}{Provide examples of legislative and regulatory reforms that were
introduced after the 2007--2009 financial crisis.}

\orsub{Reforms after the global financial crisis}

\begin{enumerate}
\item[\textcolor{rust}{\sffamily\bfseries 1)}] Creation of the \textbf{Financial
Stability Oversight Council (FSOC)} to identify and address risks threatening the
entire financial system.
\item[\textcolor{rust}{\sffamily\bfseries 2)}] \textbf{Systemically important
financial institutions (SIFIs)} may be required to hold additional capital to act
as a buffer during emergencies.
\item[\textcolor{rust}{\sffamily\bfseries 3)}] Increased oversight of
over-the-counter (OTC) derivatives, requiring clearing through exchanges for more
transparency.
\item[\textcolor{rust}{\sffamily\bfseries 4)}] Establishment of the
\textbf{Consumer Financial Protection Bureau (CFPB)} to ensure clear and fair
lending practices for consumers.
\item[\textcolor{rust}{\sffamily\bfseries 5)}] \textbf{Volcker Rule:} limits on
proprietary trading by banks that accept insured deposits.
\end{enumerate}
